\uuid{WPyj}
\titre{Bijection}
\auteur{bodin ; hanani ; mzari}
\organisation{exo7}
\licence{CC BY-NC-SA 4.0}
\datecreate{2020-01-20}
\niveau{L1}
\module{Analyse}
\chapitre{Continuité, limite et étude de fonctions réelles}
\sousChapitre{Etude de fonctions}
\difficulte{1}
\theme{Maximum, bijection, bijection}
\exo7qcmid{lille-1.continuite.026}

\contenu{
\begin{question}
\qclassification{123.04}
\qtags{motcle=bijection}
Soit $f : \Rr \to \Rr$ une fonction continue telle que : elle est strictement croissante sur $]-\infty,0]$ ; strictement décroissante sur $[0,1]$ ; strictement croissante sur $[1,+\infty[$. En plus $\lim_{x\to-\infty} f = - \infty$, $f(0)=2$, $f(1) = 1$ et $\lim_{x\to+\infty} f = 3$. Quelles sont les assertions vraies ?
\begin{answers}

    \good{La restriction $f_| : ]-\infty,0] \to ]-\infty,2]$ est bijective.}

    \bad{La restriction $f_| : [1,+\infty[ \to [1,+\infty[$ est bijective.}

    \good{La restriction $f_| : [0,1] \to [1,2]$ est bijective.}

    \bad{La restriction $f_| : ]0,+\infty] \to [1,3[$ est bijective.}
\end{answers}
\begin{explanations}
Le plus simple est de dessiner l'allure du graphe (ou le tableau de variation) pour se convaincre que $f$ restreinte à $]-\infty,0]$ définit une bijection vers $]-\infty,2]$ ;  $f$ restreinte à $[0,1]$ définit une bijection (décroissante) vers $[1,2]$ ;  $f$ restreinte à $[1,+\infty[$ définit une bijection vers $[1,3[$.
\end{explanations}
\end{question}
}
